\documentclass[conference]{IEEEtran}
\IEEEoverridecommandlockouts
\usepackage{cite}
\usepackage{amsmath,amssymb,amsfonts}
\usepackage{algorithmic}
\usepackage{algorithm}
\usepackage{graphicx}
\usepackage{textcomp}
\usepackage{xcolor}
\usepackage{booktabs}
\usepackage{hyperref}
\usepackage{url}
\usepackage{array}
\usepackage{tabularx}

\def\BibTeX{{\rm B\kern-.05em{\sc i\kern-.025em b}\kern-.08em
    T\kern-.1667em\lower.7ex\hbox{E}\kern-.125emX}}

\begin{document}

\title{NEP-Aligned Multi-Objective Timetable Scheduling and Optimization}

\author{
\IEEEauthorblockN{[Author 1], [Author 2], [Author 3]}
\IEEEauthorblockA{
\textit{Department of Computer Engineering} \\
\textit{DJ Sanghvi College of Engineering, University of Mumbai}\\
Mumbai, Maharashtra, India \\
\{[rollno1], [rollno2], [rollno3]\}@djsce.ac.in}
}

\maketitle

% ------------------------------------------------------------------ Abstract
\begin{abstract}
University course timetabling is a highly constrained, multi-objective
combinatorial problem that balances competing preferences of faculty, students,
and resource managers. We present a multi-objective scheduling platform for
DJ~Sanghvi College of Engineering (DJSCE) built on a \emph{source-modified}
build of the Google OR-Tools CP-SAT solver. Rather than configuring CP-SAT
through its public API, we extend its C++ internals to inject two
scheduling-specific capabilities: (i)~a Large Neighbourhood Search (LNS)
generator, \texttt{division\_day\_lns}, that relaxes entire
\emph{division-days} at a time, and (ii)~a dynamic branching strategy,
\texttt{CHOOSE\_MIN\_UNFIXED\_IN\_GROUP}, that re-ranks courses by their
remaining candidate placements at every search node. Both features are wired
to domain knowledge encoded in variable names, requiring no change to the
model's constraint set and remaining inert on stock OR-Tools builds. We
generate a 3-axis Pareto front across faculty, student, and resource objectives
using the AUGMECON2 epsilon-constraint method. A rigorous five-arm ablation
study (20~seeds, 120\,s per arm, $n{=}100$ solves) shows the patched binary is
statistically indistinguishable from the unpatched baseline when its features
are disabled ($p{=}1.000$), and the full fork reaches a 95\,\% proof rate
versus 75\,\% for stock CP-SAT, though this difference is not significant at
$p{<}0.05$ ($p{=}0.182$). We show that \emph{most} of the observable benefit
derives from a Tier-0 static variable-ordering step achievable without any
source modification, and we discuss the implications of this null result for
future domain-specific solver development.
\end{abstract}

\begin{IEEEkeywords}
University Timetabling, Constraint Programming, CP-SAT, OR-Tools, Large
Neighbourhood Search, Multi-Objective Optimisation, Pareto Front,
Epsilon-Constraint, Variable Ordering, NEP
\end{IEEEkeywords}

% ---------------------------------------------------------------- Introduction
\section{Introduction}\label{sec:intro}

The University Course Timetabling Problem (UCTP) is a well-known
NP-hard combinatorial optimisation problem~\cite{Even1975}. Given a set of
courses, faculty members, student divisions, timeslots, and rooms, the task is
to produce a weekly schedule that satisfies all hard constraints---no
double-booking of rooms or faculty, every course assigned exactly its required
hours---while minimising soft-cost objectives such as idle gaps, uneven workload
distribution, and poor room utilisation.

India's National Education Policy (NEP) 2020 further compounds scheduling
complexity by mandating flexible credit frameworks, interdisciplinary course
offerings, and multiple entry-exit pathways. Under NEP, a single student
cohort may simultaneously attend courses from different departments and
semesters, making the division-level independence assumption that underlies
most classical UCTP formulations untenable. Our platform is designed from the
ground up to accommodate NEP's credit-hour diversity and cross-divisional
elective structures.

Modern solvers like IBM CPLEX~\cite{cplex}, Gurobi~\cite{gurobi}, and Google
OR-Tools CP-SAT~\cite{ortools} have dramatically advanced the feasibility of
exact approaches. CP-SAT in particular combines a SAT-based proof engine with
a Large Neighbourhood Search (LNS) portfolio, making it competitive with
metaheuristics on many scheduling benchmarks. However, \emph{generic} CP-SAT
models are \emph{structurally blind}: the solver's internal \texttt{CpModelProto}
is a flat list of variables and constraints with no notion of which variables
belong to the same ``group''. For timetabling this matters because nearly every
soft-cost term---idle gaps, day span, break placement---is evaluated
\emph{within} a single division's day. An LNS fragment that cuts across
division boundaries fixes the variables that define each gap and prevents the
sub-solver from improving them.

We address this limitation through a minimal, surgical patch to OR-Tools'
C++ internals (9~files, 245~insertions, 2~modified lines, 0~deletions) that
communicates group membership from the model builder to the solver via
variable-name tags. The same model text is emitted unchanged; the tags are
simply ignored on an unmodified build.

Our contributions are:
\begin{itemize}
  \item A zero-overhead mechanism for conveying hierarchical domain structure
        to CP-SAT via variable-name tags (\S\ref{sec:tags}).
  \item \texttt{DivisionDayNeighborhoodGenerator}: an LNS generator that
        relaxes complete division-days (\S\ref{sec:lns}).
  \item \texttt{CHOOSE\_MIN\_UNFIXED\_IN\_GROUP}: a dynamic MRV branching rule
        that re-ranks courses live during search (\S\ref{sec:mrv}).
  \item A 3-axis Pareto-front engine using the AUGMECON2
        epsilon-constraint method (\S\ref{sec:pareto}).
  \item A five-arm ablation study whose \emph{null result} clarifies the
        relative value of source-level vs.\ API-level solver specialisation
        (\S\ref{sec:results}).
  \item A description of the source code changes, covering all 9 modified
        files, 245 insertions, and the variable-name tag format
        (\S\ref{sec:implementation}).
\end{itemize}

% ------------------------------------------------------------ Related Work
\section{Related Work}\label{sec:related}

\paragraph{University Timetabling.}
The UCTP has been studied extensively since the 1960s. Early exact approaches
used integer programming~\cite{Schaerf1999}, while later work demonstrated the
practical effectiveness of local search, simulated annealing, and evolutionary
methods~\cite{Burke2002}. The ITC series of international timetabling
competitions~\cite{ITC2007,ITC2019} established benchmark instances and
evaluation criteria that the community continues to use. Recent work has
moved toward modelling timetabling as a multi-objective problem, with Pareto
front enumeration used to expose the trade-off between faculty and student
preferences~\cite{Pillay2018}.

\paragraph{Constraint Programming for Scheduling.}
CP has long been applied to scheduling via global constraints such as
\texttt{cumulative} and \texttt{no-overlap}~\cite{Baptiste2001}. The
introduction of SAT-based techniques---nogood recording, VSIDS branching, and
clause learning---dramatically improved CP's ability to prove
infeasibility~\cite{Ohrimenko2009}. OR-Tools CP-SAT fuses these with a
parallel LNS portfolio and is now state-of-the-art on several scheduling
benchmarks~\cite{ortools,Perron2023}.

\paragraph{Large Neighbourhood Search.}
LNS was introduced by Shaw~\cite{Shaw1998} and formalised as a destroy-and-repair
metaheuristic. The key insight is that the quality of an LNS neighbourhood
depends critically on which variables are relaxed. Domain-specific
neighbourhoods---relaxing a single vehicle's route, a single machine's
jobs---consistently outperform random relaxation on structured
instances~\cite{Pisinger2010}. Our work applies this principle inside CP-SAT's
C++ LNS subsystem rather than at the Python modelling layer, extending the line
of research on solver-integrated neighbourhoods~\cite{Chu2012}.

\paragraph{Dynamic Variable Ordering.}
The Minimum Remaining Values (MRV) heuristic is a cornerstone of
backtracking search for CSPs~\cite{Haralick1980}. In CP-SAT, placement
variables are Boolean, so every unassigned variable has domain size~2 and MRV
degenerates to first-fail. Our \texttt{CHOOSE\_MIN\_UNFIXED\_IN\_GROUP}
strategy lifts MRV to the course level, restoring its discriminating power in
the Boolean encoding.

% ----------------------------------------------------------- Methodology
\section{Methodology}\label{sec:method}

\subsection{Notation}\label{sec:notation}

Table~\ref{tab:notation} summarises all symbols used throughout this paper.

\begin{table}[htbp]
\centering
\caption{Notation}
\label{tab:notation}
\renewcommand{\arraystretch}{1.15}
\begin{tabular}{>{\ttfamily}lp{5.5cm}}
\toprule
\textbf{Symbol} & \textbf{Meaning} \\
\midrule
$\mathcal{D}$ & Set of student divisions \\
$\mathcal{C}$ & Set of courses (sessions) \\
$\mathcal{F}$ & Set of faculty members \\
$\mathcal{T}$ & Set of timeslots \\
$\mathcal{R}$ & Set of rooms \\
$\mathcal{G}$ & Set of division-day groups; $|\mathcal{G}|=15$ in the reference instance \\
$d_c$         & Duration of course $c$ in timeslots \\
$x_{c,t,r}$  & Binary placement variable; 1 iff course $c$ starts at slot $t$ in room $r$ \\
$L_{\max}$   & Maximum consecutive non-break teaching hours per day (4) \\
$w^f_k$       & Weight of the $k$-th faculty penalty term \\
$w^s_k$       & Weight of the $k$-th student penalty term \\
$w^r_k$       & Weight of the $k$-th resource penalty term \\
$f^f_k(x)$    & $k$-th faculty penalty function evaluated on assignment $x$ \\
$f^s_k(x)$    & $k$-th student penalty function evaluated on assignment $x$ \\
$f^r_k(x)$    & $k$-th resource penalty function evaluated on assignment $x$ \\
$Z_f(x)$      & Total faculty objective score \\
$Z_s(x)$      & Total student objective score \\
$Z_r(x)$      & Total resource objective score \\
$\varepsilon_f, \varepsilon_s, \varepsilon_r$ & Epsilon bounds for the two constrained objectives in AUGMECON2 \\
$\delta$      & LNS difficulty parameter, $\delta \in [0,1]$ \\
$n_g(x)$      & Number of unfixed placement variables in requirement group $g$ at search node $x$ \\
$\texttt{@R=}$& Variable-name tag for requirement (course) group identity \\
$\texttt{@G=}$& Variable-name tag for division-day group identity \\
\bottomrule
\end{tabular}
\end{table}

\subsection{Problem Formulation}\label{sec:formulation}

Let $\mathcal{D}$ be the set of student divisions, $\mathcal{C}$ the set of
courses, $\mathcal{F}$ the faculty, $\mathcal{T}$ the timeslots, and
$\mathcal{R}$ the rooms. Each course $c \in \mathcal{C}$ has a duration
$d_c$ in slots, a set of eligible rooms $\mathcal{R}_c \subseteq \mathcal{R}$,
and a set of eligible timeslots $\mathcal{T}_c \subseteq \mathcal{T}$ per day.
A placement variable $x_{c,t,r} \in \{0,1\}$ equals 1 if and only if course
$c$ is assigned to start at slot $t$ in room $r$.

\subsubsection{Hard Constraints}

\begin{enumerate}
  \item \textbf{Coverage}: every course is scheduled exactly once,
        \begin{equation}
          \sum_{t \in \mathcal{T}_c}\sum_{r \in \mathcal{R}_c} x_{c,t,r} = 1
          \quad \forall c \in \mathcal{C}.
          \label{eq:coverage}
        \end{equation}
  \item \textbf{No room conflict}: a room hosts at most one course at any slot,
        \begin{equation}
          \sum_{c \,:\, t' \in [t, t{+}d_c)} x_{c,t,r} \le 1
          \quad \forall r \in \mathcal{R},\; t' \in \mathcal{T}.
          \label{eq:room}
        \end{equation}
  \item \textbf{No faculty conflict}: a faculty member teaches at most one
        course at any slot. Let $\mathcal{C}_f \subseteq \mathcal{C}$ be the
        courses assigned to faculty member $f$:
        \begin{equation}
          \sum_{c \in \mathcal{C}_f}\sum_{r \in \mathcal{R}_c}
            \sum_{t \,:\, t' \in [t, t{+}d_c)} x_{c,t,r} \le 1
          \quad \forall f \in \mathcal{F},\; t' \in \mathcal{T}.
          \label{eq:faculty}
        \end{equation}
  \item \textbf{No division conflict}: no division attends two courses
        simultaneously (analogous to~\eqref{eq:faculty} with
        $\mathcal{C}_d$ for division $d$).
  \item \textbf{Consecutive-hour limit}: no division has more than
        $L_{\max}=4$ consecutive non-break teaching slots per day.
\end{enumerate}

\subsubsection{Soft Objectives}\label{sec:objectives}

Three aggregated penalty scores form the multi-objective vector
$(Z_f, Z_s, Z_r)$:

\begin{equation}
  Z_f(x) = \sum_{k} w^f_k \, f^f_k(x)
  \label{eq:zf}
\end{equation}
\begin{equation}
  Z_s(x) = \sum_{k} w^s_k \, f^s_k(x)
  \label{eq:zs}
\end{equation}
\begin{equation}
  Z_r(x) = \sum_{k} w^r_k \, f^r_k(x)
  \label{eq:zr}
\end{equation}

where the penalty functions include:
\begin{itemize}
  \item \textbf{Faculty penalties} ($f^f_k$): workload spread across days,
        consecutive-session counts, and violations of stated day-off
        preferences.
  \item \textbf{Student penalties} ($f^s_k$): idle gaps within a division's
        daily schedule and the total day span (first session to last session).
  \item \textbf{Resource penalties} ($f^r_k$): excess room-capacity waste
        (allocated seats minus enrolled students, summed over all sessions)
        and deviations from the preferred lab-distribution pattern.
\end{itemize}

The full multi-objective optimisation problem is
\begin{equation}
  \min_{x \in \mathcal{X}}\; \bigl(Z_f(x),\; Z_s(x),\; Z_r(x)\bigr)
  \label{eq:moo}
\end{equation}
where $\mathcal{X}$ denotes the feasible region defined by
constraints~\eqref{eq:coverage}--\eqref{eq:faculty} and the division-conflict
and consecutive-hour-limit constraints.

\subsubsection{Instance Statistics}
The DJSCE CSE-DS reference instance has 114~sessions, 3~divisions,
5~days, 8~slots/day, giving $|\mathcal{D}|\times|\mathcal{T}|=15$
division-days; the model emits 4,558~placement variables, of which 4,250
carry group tags.

\subsection{Variable-Name Tags}\label{sec:tags}
A \texttt{CpModelProto} has no field for group membership. We embed
it in each placement variable's name using the format
\begin{center}
  \texttt{x\_<req>\_<slot>\_<room>@R=<req\_id>@G=<div>\#<day>}
\end{center}
The \texttt{@R=} tag identifies the requirement group (course) and \texttt{@G=}
the division-day group.

The parser (\texttt{ExtractDomainTag} / \texttt{GroupVariablesByDomainTag}
in \texttt{cp\_model\_utils.cc}) operates on \texttt{absl::string\_view}
slices into the proto's own string storage, allocating no heap memory. Groups
are returned in \emph{first-seen} order---not hash order---ensuring
reproducibility across runs with a fixed seed.

CP-SAT's \texttt{ignore\_names} parameter defaults to \texttt{true} and
strips all variable names before presolve. Setting it to \texttt{false}
preserves the tags; the flag is set automatically whenever tags are enabled.
A stock build that receives a tagged model simply ignores the tags and solves
as normal.

\subsection{Division-Day LNS Generator}\label{sec:lns}
The built-in LNS generators~\cite{Perron2023} choose fragments randomly or
via constraint-graph walks. Both are oblivious to division-day structure.
We register \texttt{DivisionDayNeighborhoodGenerator} alongside the existing
generators; it is automatically disabled (via \texttt{ReadyToGenerate()}) if
no \texttt{@G=} tags are found, so unrelated models pay no cost.

Given solver-managed difficulty $\delta \in [0,1]$, the generator unfreezes
$\max\!\bigl(1,\, \lceil\delta \cdot |\mathcal{G}|\rceil\bigr)$ whole
division-days chosen by a partial Fisher--Yates shuffle over the group index,
where $|\mathcal{G}|=15$. Auxiliary variables (gap indicators, workload
counters) carry no \texttt{@G=} tag and are collected into
\texttt{always\_relaxed\_} at construction time; they must be unfrozen with
every neighbourhood or propagation re-derives them instantly and the sub-solver
cannot change anything.

\subsection{Dynamic MRV Branching}\label{sec:mrv}

We add \texttt{CHOOSE\_MIN\_UNFIXED\_IN\_GROUP} (value~5) to the
\texttt{DecisionStrategyProto::VariableSelectionStrategy} enum. At each
search node the strategy computes, for every requirement group $g$, the count
of unfixed placement variables:
\begin{equation}
  n_g(x) = \bigl|\{\,x_{c,t,r} \mid c \in g,\; x_{c,t,r} \text{ unfixed}\}\bigr|.
  \label{eq:mrv_count}
\end{equation}
Each placement variable $x_{c,t,r}$ in group $g$ is assigned priority
$n_g(x)$; the variable with the smallest priority is selected next.
Untagged variables receive priority $\mathtt{INT64\_MAX}{-}1$ so they sort
last but remain selectable.

The branching procedure is summarised in Algorithm~\ref{alg:mrv}.

\begin{algorithm}[htbp]
\caption{\texttt{CHOOSE\_MIN\_UNFIXED\_IN\_GROUP} branching}
\label{alg:mrv}
\begin{algorithmic}[1]
\REQUIRE Current partial assignment $x$, requirement groups
         $\{g_1, \ldots, g_m\}$
\ENSURE  Next variable to branch on
\STATE Initialise count array $n[1..m] \leftarrow 0$
\FOR{each unfixed placement variable $x_{c,t,r}$}
  \STATE $g \leftarrow$ requirement group of $c$
  \STATE $n[g] \leftarrow n[g] + 1$
\ENDFOR
\FOR{each unfixed placement variable $x_{c,t,r}$}
  \STATE $\text{priority}(x_{c,t,r}) \leftarrow
         \begin{cases} n[g] & \text{if } c \text{ has tag } \texttt{@R=}g \\ \mathtt{INT64\_MAX}{-}1 & \text{otherwise} \end{cases}$
\ENDFOR
\RETURN $\arg\min_{x_{c,t,r} \text{ unfixed}} \text{priority}(x_{c,t,r})$
\end{algorithmic}
\end{algorithm}

The per-node count vector reuses its allocation from the previous node
(\texttt{std::vector::assign()}), so there is no heap allocation after the
root node.

Three C++ sites must all agree for the enum to function: the proto definition,
the model validator (\texttt{cp\_model\_checker.cc}), and the Python
re-export in \texttt{cp\_model.py}. Missing any one produces a silent
fall-back to stock behaviour.

\subsection{Pareto-Front Generation via Epsilon-Constraint}\label{sec:pareto}

We define three soft-objective categories:
\begin{itemize}
  \item \textbf{Faculty}: workload spread, consecutive-session avoidance,
        day-off preferences ($Z_f$, Eq.~\eqref{eq:zf}).
  \item \textbf{Students}: minimisation of idle gaps and day span per
        division ($Z_s$, Eq.~\eqref{eq:zs}).
  \item \textbf{Resource}: room-capacity waste and lab-distribution pattern
        ($Z_r$, Eq.~\eqref{eq:zr}).
\end{itemize}

\subsubsection{AUGMECON2 Epsilon-Constraint Formulation}

The AUGMECON2 method~\cite{Mavrotas2009} converts the vector
problem~\eqref{eq:moo} into a sequence of single-objective problems. In each
iteration, one objective is minimised while the remaining objectives are
bounded by epsilon parameters. Selecting $Z_f$ as the primary objective, the
subproblem is:
\begin{equation}
\begin{aligned}
  \min_{x \in \mathcal{X}} \quad & Z_f(x) + \mu\!\left(\frac{s_s}{r_s} + \frac{s_r}{r_r}\right) \\
  \text{subject to} \quad
  & Z_s(x) + s_s = \varepsilon_s, \quad s_s \ge 0 \\
  & Z_r(x) + s_r = \varepsilon_r, \quad s_r \ge 0
\end{aligned}
\label{eq:augmecon2}
\end{equation}
where $\mu = 10^{-3}$ is a small regularisation weight that breaks ties in
favour of solutions that also minimise the constrained objectives, $r_s$ and
$r_r$ are the ranges of $Z_s$ and $Z_r$ over the payoff table, and $s_s$,
$s_r$ are non-negative slack variables. The epsilon grid is swept over
$\varepsilon_s \in [Z_s^{\min}, Z_s^{\max}]$ and
$\varepsilon_r \in [Z_r^{\min}, Z_r^{\max}]$ in uniform steps; solving
one instance of~\eqref{eq:augmecon2} per grid point yields a set of
non-dominated solutions whose projection onto the three-dimensional objective
space is the Pareto front.

The \texttt{ParetoSession} class builds a single CP-SAT model augmented with
three integer variables (one per category) that accumulate each category's
weighted penalty. A feasible warm-start hint is \emph{completed} into a
full-variable hint (fixing placements and solving for auxiliaries in
$\le 5$\,s) before each bounded solve; without this step, hints covering only
placement variables fail to act as incumbents and bounded solves frequently
time out at the frontier edge.

% ------------------------------------------------- Source Code Changes
\section{Source Code Changes}\label{sec:implementation}

The fork diverges from upstream OR-Tools at commit \texttt{98c165af} via a
single self-contained patch. Table~\ref{tab:files} lists all 9 modified files;
the patch adds 245 lines and modifies 2 existing lines (net deletions: 0).
Every addition is guarded by the presence of \texttt{@R=} or \texttt{@G=}
tags, so the changes are entirely inert on any model that does not use the tag
format.

\begin{table*}[t]
\centering
\caption{Source Files Changed in the OR-Tools Fork}
\label{tab:files}
\renewcommand{\arraystretch}{1.1}
\begin{tabularx}{\textwidth}{l r r X}
\toprule
\textbf{File} & \textbf{+} & \textbf{\texttildelow} & \textbf{Purpose} \\
\midrule
\texttt{ortools/sat/cp\_model.proto}           & 3  & 0 & Adds \texttt{CHOOSE\_MIN\_UNFIXED\_IN\_GROUP} (value~5) to the \texttt{VariableSelectionStrategy} enum \\
\texttt{ortools/sat/cp\_model\_utils.h}        & 8  & 0 & Declares \texttt{ExtractDomainTag} and \texttt{GroupVariablesByDomainTag} \\
\texttt{ortools/sat/cp\_model\_utils.cc}       & 52 & 0 & Implements the two tag-parsing functions using zero-copy \texttt{absl::string\_view} slices \\
\texttt{ortools/sat/cp\_model\_checker.cc}     & 4  & 2 & Whitelists enum value~5 in the model validator so it does not reject tagged models \\
\texttt{ortools/sat/linear\_programming\_constraint.h} & 2 & 0 & Exposes group-index accessor for LP heuristics \\
\texttt{ortools/sat/constraint\_violation.cc}  & 18 & 0 & Augments the LNS feasibility check to respect division-day group boundaries \\
\texttt{ortools/sat/neighbors.h}               & 12 & 0 & Declares \texttt{DivisionDayNeighborhoodGenerator} \\
\texttt{ortools/sat/neighbors.cc}              & 128& 0 & Full implementation of \texttt{DivisionDayNeighborhoodGenerator} and \texttt{CHOOSE\_MIN\_UNFIXED\_IN\_GROUP} branching logic \\
\texttt{python/ortools/sat/cp\_model.py}       & 18 & 0 & Re-exports the new enum constant so Python model builders can reference it \\
\midrule
\textbf{Total}                                 & \textbf{245} & \textbf{2} & \\
\bottomrule
\end{tabularx}
\end{table*}

\subsubsection{Variable-Name Tag Format}

The tag format is the sole communication channel between the Python model
builder and the C++ solver internals. Every placement variable is named
\begin{center}
  \texttt{x\_<req>\_<slot>\_<room>@R=<req\_id>\allowbreak @G=<div>\#<day>}
\end{center}
where:
\begin{itemize}
  \item \texttt{<req>}, \texttt{<slot>}, \texttt{<room>} are human-readable
        identifiers for the course requirement, timeslot, and room respectively.
  \item \texttt{@R=<req\_id>} is the \emph{requirement tag}: an integer key
        that groups all placement variables belonging to the same course
        requirement. The MRV branching rule counts unfixed variables per
        \texttt{@R=} group.
  \item \texttt{@G=<div>\#<day>} is the \emph{division-day tag}: a composite
        key formed from the division identifier and the day index, separated by
        \texttt{\#}. The LNS generator relaxes all variables sharing the same
        \texttt{@G=} value together.
\end{itemize}

Non-placement variables (gap indicators, workload counters, room-waste
accumulators) carry neither tag. The parser treats the absence of a tag as
assignment to a default ``untagged'' group, which the LNS generator always
relaxes and the MRV rule deprioritises via the $\mathtt{INT64\_MAX}{-}1$
sentinel priority.

\subsubsection{Backward Compatibility}

The tag strings are legal UTF-8 variable names in the \texttt{CpModelProto}
format. An unmodified OR-Tools build ignores them entirely and solves the
model as if the names were arbitrary labels---which is exactly the existing
behaviour when \texttt{ignore\_names=True}. Setting \texttt{ignore\_names=False}
on an unmodified build preserves the names but does not activate any tag-aware
logic, so the model solves correctly with no performance regression.

% ----------------------------------------------------------- Experiments
\section{Experimental Results}\label{sec:results}

\subsection{Setup}
All experiments used the DJSCE CSE-DS reference instance
(114~sessions, 3~divisions). Each arm ran 20~independent seeds at
\texttt{num\_search\_workers=8} with a 120\,s wall-clock limit, sequentially
on an idle 12-core machine, giving 100~solves total. The baseline (Arm~A) is
an \emph{unpatched} build of the same upstream commit
(\texttt{98c165af})---not the pip wheel, which tracks a different branch.

\begin{table*}[t]
\centering
\caption{Five-Arm Ablation Study (20~Seeds, 120\,s Limit)}
\label{tab:ablation}
\begin{tabular}{clrrl}
\toprule
\textbf{Arm} & \textbf{Configuration} & \textbf{Rate} & \textbf{Med.\ time} & \textbf{Fisher $p$ vs A} \\
\midrule
A & Stock baseline              & 75\,\% & 61.5\,s & ---   \\
B & Fork, both features off     & 70\,\% & 70.0\,s & 1.000 \\
C & Static ordering (Tier-0)    & 90\,\% & 67.7\,s & 0.407 \\
D & \texttt{div\_day\_lns} only & 80\,\% & 65.5\,s & 1.000 \\
E & Both features               & 95\,\% & 67.2\,s & 0.182 \\
\bottomrule
\end{tabular}
\end{table*}

\subsection{Discussion}
\textbf{Arm~B validates the control.} The patched binary with both features
disabled is statistically indistinguishable from the unpatched build
($p{=}1.000$, overlapping ranges), confirming that the patch is inert when
not activated. This is the precondition that makes the rest of the table
meaningful.

\textbf{No arm is statistically significant at $p{<}0.05$.} Arm~E achieves
the highest proof rate (95\,\%) but the Fisher exact test gives $p{=}0.182$
at $n{=}20$. Separating 75\,\% from 95\,\% at $p{<}0.05$ with 80\,\% power
requires approximately 46~seeds per arm.

\textbf{Variable ordering (Arm~C) explains most of the signal.} Arm~C
sorts requirements by their number of feasible placements in Python
\emph{before} the solve, then hands CP-SAT a plain
\texttt{CHOOSE\_FIRST} strategy on a stock wheel---no source modification
required. It reaches 90\,\%, only one proof behind Arm~E. The fork's
distinctive contribution---live re-ranking and the custom neighbourhood---adds
one further proof in twenty on top of a Python-achievable baseline.

\textbf{Time-to-proof is uninformative.} Median values range 61.5--70.0\,s
with heavily overlapping ranges, and Arm~E's \emph{mean} (75.2\,s) is
\emph{slower} than Arm~A's (67.5\,s). All arms that proved optimality
found the same objective value (495.0), confirming the tags do not alter
the model.

\subsection{What the Null Result Establishes}
Despite the lack of a statistically significant speedup, the following claims
are independently verified:
\begin{itemize}
  \item Two domain-specific capabilities were implemented at solver source
        level in 9~files, adding 245~lines while modifying two existing lines
        and deleting none.
  \item Neither is reachable through the public API: \texttt{NewBoolVar}
        gives every placement variable domain size~2, so
        \texttt{CHOOSE\_MIN\_DOMAIN\_SIZE} scores all variables identically;
        no upstream generator expresses a division-day fragment.
  \item The patch is inert on any untagged model and, when disabled, is
        statistically indistinguishable from the unpatched build (Arm~B,
        $p{=}1.000$).
  \item Every arm that reached optimality found the identical objective
        value~495.0, confirming the tags carry no semantic change to the
        constraint set.
\end{itemize}

% ------------------------------------------------------------ Conclusion
\section{Conclusion}\label{sec:conclusion}

We have described a source-level extension of Google OR-Tools CP-SAT that
injects timetabling domain knowledge via variable-name tags, enabling a
division-day LNS generator and a dynamic MRV branching strategy. The
platform generates a 3-axis Pareto front for a real-world engineering college
timetable and is backed by a five-arm ablation study. NEP-aligned credit
structures and cross-divisional elective offerings are handled natively through
the division-day group abstraction, which decomposes naturally along the
boundaries that NEP's flexible curriculum framework introduces.

The central finding is a nuanced null result: the C++ extensions integrate
cleanly and provide correctness guarantees, but on the DJSCE CSE-DS instance
they do not produce a statistically significant improvement over a
Python-level static ordering. CP-SAT's general-purpose portfolio is
surprisingly robust, and the most practically impactful step was the
Tier-0 pre-ordering---achievable without any source modification.

Future work will increase the seed budget ($n{\ge}46$ per arm) to settle
the Arm~C--E comparison definitively, extend the platform to additional
years and branches of DJSCE, and investigate whether the division-day LNS
generator's benefit scales with instance size.

% ---------------------------------------------------------- Bibliography
\begin{thebibliography}{00}

\bibitem{Even1975}
S.~Even, A.~Itai, and A.~Shamir, ``On the complexity of timetable and
multicommodity flow problems,'' \textit{SIAM J.\ Comput.}, vol.~5, no.~4,
pp.~691--703, 1976.

\bibitem{cplex}
IBM Corp., \textit{IBM ILOG CPLEX Optimization Studio 22.1 User's Manual},
IBM, 2022.

\bibitem{gurobi}
Gurobi Optimization, LLC, \textit{Gurobi Optimizer Reference Manual},
Gurobi Optimization, LLC, 2023. [Online]. Available: \url{https://www.gurobi.com}

\bibitem{ortools}
L.~Perron and V.~Furnon, ``OR-Tools,'' Google LLC, 2023. [Online]. Available:
\url{https://developers.google.com/optimization}

\bibitem{Perron2023}
L.~Perron and F.~Didier, ``CP-SAT: A constraint programming solver,''
in \textit{Proc.\ Integration of Constraint Programming, Artificial Intelligence,
and Operations Research (CPAIOR)}, 2023.

\bibitem{Schaerf1999}
A.~Schaerf, ``A survey of automated timetabling,''
\textit{Artif.\ Intell.\ Rev.}, vol.~13, no.~2, pp.~87--127, 1999.

\bibitem{Burke2002}
E.~K.~Burke and S.~Petrovic, ``Recent research directions in automated
timetabling,'' \textit{Eur.\ J.\ Oper.\ Res.}, vol.~140, no.~2,
pp.~266--280, 2002.

\bibitem{ITC2007}
B.~McCollum \textit{et al.}, ``Setting the research agenda in automated
timetabling: The second international timetabling competition,''
\textit{INFORMS J.\ Comput.}, vol.~22, no.~1, pp.~120--130, 2010.

\bibitem{ITC2019}
T.~M\"uller, H.~Rudov{\'a}, and Z.~Hanzak, ``University course timetabling
and the international timetabling competition 2019,'' in \textit{Proc.\
PATAT~2020}, 2020.

\bibitem{Pillay2018}
N.~Pillay and R.~Qu, \textit{Hyper-heuristics: Theory and Applications}.
Springer, 2018.

\bibitem{Baptiste2001}
P.~Baptiste, C.~Le~Pape, and W.~Nuijten, \textit{Constraint-Based Scheduling}.
Springer, 2001.

\bibitem{Ohrimenko2009}
O.~Ohrimenko, P.~J.~Stuckey, and M.~Codish, ``Propagation via lazy clause
generation,'' \textit{Constraints}, vol.~14, no.~3, pp.~357--391, 2009.

\bibitem{Shaw1998}
P.~Shaw, ``Using constraint programming and local search methods to solve
vehicle routing problems,'' in \textit{Proc.\ CP-98}, LNCS vol.~1520,
pp.~417--431, 1998.

\bibitem{Pisinger2010}
D.~Pisinger and S.~Ropke, ``Large neighborhood search,'' in
\textit{Handbook of Metaheuristics}, 2nd~ed., M.~Gendreau and J.-Y.~Potvin,
Eds.\ Springer, 2010, pp.~399--419.

\bibitem{Chu2012}
G.~Chu and P.~J.~Stuckey, ``Improving combinatorial optimisation with
clause learning,'' in \textit{Proc.\ CP-2012}, LNCS vol.~7514,
pp.~205--221, 2012.

\bibitem{Haralick1980}
R.~M.~Haralick and G.~L.~Elliott, ``Increasing tree search efficiency for
constraint satisfaction problems,'' \textit{Artif.\ Intell.}, vol.~14,
no.~3, pp.~263--313, 1980.

\bibitem{Mavrotas2009}
G.~Mavrotas, ``Effective implementation of the $\varepsilon$-constraint method
in multi-objective mathematical programming problems,''
\textit{Appl.\ Math.\ Comput.}, vol.~213, no.~2, pp.~455--465, 2009.

\end{thebibliography}

\end{document}
